% 原书物理页=印刷页8--12上；1.3 全节（起于 p8 中部节题，止于 p12 上部第 2 章章名之前）.
\section{Lie 群的 Lie 代数}\label{sec:ssl-1-3}

有多种方式在 Lie 群的单位元处的切空间 $T_eG$ 上赋予一个 $\mathbb{R}$-Lie 代数结构.

\begin{numpara}[经由左不变向量场]\label{np:ssl-1-3-1}
这可能是最快的一种做法.想法是：$e$ 处的每个切向量，借助\newterm{ssl-left-translation}{左平移}{left translation}作用确定一个\newterm{ssl-left-invariant-vector-field}{左不变向量场}{left-invariant vector field}，并且这一构造是双射.形式地，让 $G$ 通过左平移作用在自身上，记 $L_g(x)=gx$.这在切丛 $TG$ 上诱导出一个与投影 $\pi$ 相容的作用，我们把它记作 $(g,\theta)\mapsto L^{t}_{g}(\theta)$.映射
\[
\begin{array}{rcl}
TG &\longrightarrow& G\times T_eG\\
\theta &\longmapsto& \bigl(\pi(\theta),\,L^{t}_{\pi^{-1}(\theta)}(\theta)\bigr)
\end{array}
\]
\footnote{译注：原文下标如此；按上下文应为 $\pi(\theta)^{-1}$（底点 $\pi(\theta)$ 的逆元）之误，如此第二个分量才落在 $T_eG$ 中.}
是与到 $G$ 上的投影相容的微分同胚，其逆由
\[
\begin{array}{rcl}
G\times T_eG &\longrightarrow& TG\\
(g,X) &\longmapsto& (g,\,d_eL_g(X)).
\end{array}
\]
给出.于是 $G$ 的切丛是``\newterm{ssl-parallelizable}{可平行化}{parallelizable}''的.取 $X\in\operatorname{Lie}(G):=T_eG$.于是可由下述公式为它对应一个左不变向量场 $\xi_X$（即对所有 $h,g\in G$ 满足 $\xi(gh)=d_h(L_g)(\xi(h))$ 的向量场）：
\[
\forall g\in G,\quad\xi_X(g):=d_eL_g(X).
\]
设 $V(G)$ 为 $G$ 上向量场构成的 $\mathbb{R}$-向量空间，$V(G)^{L}$ 为其中左不变向量场所成的子空间.\textbf{显然映射 $X\mapsto\xi_X$ 是 $\mathbb{R}$-线性同构 $T_eG\xrightarrow{\ \sim\ }V(G)^{L}$，其逆为 $\xi\mapsto\xi(e)$.}现在，与左不变向量场相伴的导子仍是左不变的，并且一个计算表明：两个左不变导子的Lie 括号仍左不变.这就在 $T_eG$ 上定义出一个 Lie 括号.
\end{numpara}

\begin{numpara}[经由 $G$ 的伴随表示的微分]\label{np:ssl-1-3-2}
对元素 $x\in G$，记
\[
\begin{array}{rcl}
\operatorname{Int}_x\colon G &\longrightarrow& G\\
y &\longmapsto& xyx^{-1}
\end{array}
\]
以及 $\operatorname{Ad}_x:=d_e(\operatorname{Int}_x)\colon T_eG\longrightarrow T_eG$.由于 $\operatorname{Int}_x\circ\operatorname{Int}_{x'}=\operatorname{Int}_{xx'}$，也有 $\operatorname{Ad}_x\circ\operatorname{Ad}_{x'}=\operatorname{Ad}_{xx'}$，于是 $\operatorname{Ad}_x$ 可逆，且映射
\[
\begin{array}{rcl}
\operatorname{Ad}\colon G &\longrightarrow& \mathrm{GL}(T_eG)\\
x &\longmapsto& \operatorname{Ad}_x
\end{array}
\]
是一个群同态.由于映射 $G\times G\longrightarrow G$、$(x,y)\mapsto xyx^{-1}$ 是光滑的，经微分得到的映射 $G\times T_eG\longrightarrow T_eG$、$(x,Y)\mapsto\operatorname{Ad}_x(Y)$ 也光滑，因此态射 $\operatorname{Ad}$ 是光滑的.我们把它称为 $G$ 的伴随表示.它的底空间就是 $e$ 处的切空间.由于 $\mathrm{GL}(T_eG)$ 是 $\mathbb{R}$-向量空间 $\operatorname{End}_{\mathbb{R}}(T_eG)$ 中的开集，可以把后者等同于 $\mathrm{GL}(T_eG)$ 在每一点处（特别地在点 $e$ 处）的切空间.在点 $e$ 处对 $\operatorname{Ad}$ 微分，就得到 $\mathbb{R}$-线性映射
\[
\begin{array}{rcl}
\operatorname{ad}:=d_e\operatorname{Ad}\colon T_eG &\longrightarrow& \operatorname{End}_{\mathbb{R}}(T_eG)\\
X &\longmapsto& \operatorname{ad}_X .
\end{array}
\]
只需再按下式定义括号：
\[
[X,Y]:=\operatorname{ad}_X(Y),\quad\text{对 }X,Y\in T_eG .
\]
按构造，它是双线性映射 $L\times L\longrightarrow L$.Lie 括号公理 i) 与 ii) 的验证，见 [GALM2] 1.5.4 与 1.5.6.至于此括号与前一个括号一致的验证，见 [GALM2] 1.8.4.

\noindent\textbf{性质.}
\begin{enumerate}[label=\roman*)]
\item 任何 Lie 群同态 $\varphi\colon G\longrightarrow G'$ 都通过在单位元处求导，诱导 Lie 代数同态 $d\varphi\colon\operatorname{Lie}(G)\longrightarrow\operatorname{Lie}(G)$\footnote{译注：原文此处作 $\operatorname{Lie}(G)$，应为 $\operatorname{Lie}(G')$ 之误.}.
\item $\operatorname{Lie}(G\times G')=\operatorname{Lie}(G)\times\operatorname{Lie}(G')$（显然）.
\item $\operatorname{Lie}(G)=\operatorname{Lie}(G^{0})$（含 $e$ 的\newterm{ssl-connected-component}{连通分支}{connected component}）.
\end{enumerate}

\begin{example*}
\begin{enumerate}[label=\roman*)]
\item 离散群的 Lie 代数为零.
\item $\operatorname{Lie}(\mathbb{R})=\operatorname{Lie}(\mathbb{R}^{\times})=\operatorname{Lie}(S^{1})=\mathbb{R}$（带有零括号）.更一般地，\textbf{Abel}~Lie 群的 Lie 括号为零.
\item $\operatorname{Lie}(\mathrm{GL}_n(\mathbb{R}))=\mathfrak{gl}_n(\mathbb{R})$（记号见前文），且伴随表示由共轭给出.事实上，在 $I_n$ 处的切空间等同于 $M_n(\mathbb{R})$.取 $x,y\in\mathrm{GL}_n(\mathbb{R})$，设 $y=I_n+Y$ 且 $\lVert Y\rVert<1$.等式 $xyx^{-1}=I_n+xYx^{-1}$ 表明 $\operatorname{Ad}_x$ 就是 $x$ 在 $M_n(\mathbb{R})$ 上的共轭.再设 $x=I_n+X$ 且 $\lVert X\rVert<1$.于是 $\operatorname{Ad}_x(Y)=(I_n+X)Y(I_n-X)+o(X)=Y+[X,Y]+o(X)$，这表明 $\operatorname{ad}_X(Y)=[X,Y]$.
\item 更内在地说，$\operatorname{Lie}(\mathrm{GL}(V))=\mathfrak{gl}(V)$.
\end{enumerate}
\end{example*}

\begin{exercise*}
证明行列式 $\det\colon\mathrm{GL}_n(\mathbb{R})\longrightarrow\mathbb{R}^{\times}$ 的微分是迹 $\operatorname{tr}\colon\mathfrak{gl}_n(\mathbb{R})\longrightarrow\mathbb{R}$.
\end{exercise*}
\end{numpara}

\begin{numpara}[指数映射]\label{np:ssl-1-3-3}
光滑流形上微分方程的理论表明，存在\textbf{唯一}的映射
\[
\exp_G\colon\operatorname{Lie}(G)\longrightarrow G
\]
使得对每个 $X$，映射 $t\in\mathbb{R}\mapsto\exp_G(tX)$ 都是 $G$ 的一个单参数子群，且其在 $0$ 处的微分等于 $X$.具体地，该单参数子群是向量场 $\xi_X$ 过 $e$ 的\newterm{ssl-integral-curve}{积分曲线}{integral curve}；或者等价地说，$\exp_G(X)$ 是向量场 $\xi_X$ 的\newterm{ssl-flow}{流}{flow}在点 $e\in G$、时刻 $t=1$ 处的取值，参见 [GALM2] 1.7.

\begin{example*}
对 $G=\mathrm{GL}_n(\mathbb{R})$，$\exp_G$ 就是通常的\newterm{ssl-exponential-map}{指数映射}{exponential map}，定义为 $\exp(X)=\sum_{n\in\mathbb{N}}X^{n}/n!$，$X\in\mathfrak{gl}_n(\mathbb{R})$.
\end{example*}

指数映射的唯一性蕴含：若 $\varphi\colon G\longrightarrow G'$ 是光滑同态，则
\[
\exp_{G'}\circ d\varphi=\varphi\circ\exp_G .
\]
特别地，Lie 子群的指数映射是所在大群的指数映射的限制.Lie 群与其 Lie 代数之间极强的相互作用源自下面这条性质，它同样可由带参数的微分方程理论导出（或者在 $\mathrm{GL}_n(\mathbb{R})$ 的情形可以直接看出）：

\noindent\textbf{映射 $\exp_G$ 是光滑的，并且实现从 $\operatorname{Lie}(G)$ 中 $0$ 的一个邻域到 $G$ 中 $e$ 的一个邻域上的微分同胚.此外，$\exp_G$ 的像生成 $G$ 的含 $e$ 的连通分支 $G^{0}$.}
\end{numpara}

\begin{numpara}[经典 Lie 代数]\label{np:ssl-1-3-4}
Cartan--von Neumann 定理告诉我们：Lie 群 $G$ 的闭子群 $H$ 是一个子流形.因此它是一个 Lie 群，其 Lie 代数 $\operatorname{Lie}(H)$ 是 $\operatorname{Lie}(G)$ 的 Lie 子代数，即其括号由 $\operatorname{Lie}(G)$ 的括号诱导.另外，指数映射给出 $\operatorname{Lie}(H)$ 作为 $\operatorname{Lie}(G)$ 的子集的如下刻画：
\[
\operatorname{Lie}(H)=\{X\in\operatorname{Lie}(G),\ \forall t\in\mathbb{R},\ \exp_G(tX)\in H\}.
\]
把它应用于 $G=\mathrm{GL}_n(\mathbb{R})$ 或 $\mathrm{GL}(V)$，就可以计算经典群的 Lie 代数.

例如，等式 $\exp(\operatorname{tr}(u))=\det(\exp(u))$ 表明
\[
\mathfrak{sl}(V):=\operatorname{Lie}(\mathrm{SL}(V))=\{u\in\mathfrak{gl}(V),\ \operatorname{tr}(u)=0\}.
\]
注意它总是一个 $K$-向量空间（而不只是 $\mathbb{R}$-向量空间）.用矩阵的语言，这给出
\begin{enumerate}[label=\roman*)]
\item $\mathfrak{sl}_n(\mathbb{R})=\{X\in\mathfrak{gl}_n(\mathbb{R}),\ \operatorname{tr}(X)=0\}$（迹为零的矩阵）.其维数为 $n^{2}-1$.
\item $\mathfrak{sl}_n(\mathbb{C})=\{X\in\mathfrak{gl}_n(\mathbb{C}),\ \operatorname{tr}(X)=0\}$.其实维数为 $2n^{2}-2$.
\end{enumerate}
另外，设 $G=\{g\in\mathrm{GL}(V),\ \Psi(gv,gw),\ \forall v,w\in V\}$\footnote{译注：原文此处作 $\Psi(gv,gw)$；按文意（参见下式 $gg^{*}=\mathrm{id}_{V}$）应为 $\Psi(gv,gw)=\Psi(v,w)$ 之误.}，由非退化双线性型 $\Psi$（对称、交错或 Hermite 型均可）定义.换句话说，若以 $g^{*}$ 表示 $g$ 关于 $\Psi$ 的\newterm{ssl-adjoint}{伴随元}{adjoint}，则 $G=\{g\in\mathrm{GL}(V),\ gg^{*}=\mathrm{id}_{V}\}$.由于 $\exp(u)^{*}=\exp(u^{*})$（取极限即得），我们得到
\[
\mathfrak{g}=\{u\in\mathfrak{gl}(V),\ u+u^{*}=0\}.
\]
另外注意它总是一个 $K$-向量空间，唯独酉情形例外：此时 $K=\mathbb{C}$，但 $\mathfrak{g}$ 只是一个 $\mathbb{R}$-向量空间.用矩阵的语言，这给出：
\begin{enumerate}[label=\roman*),start=3]
\item $\mathfrak{so}_n(\mathbb{R})=\operatorname{Lie}(\mathrm{O}(n))=\operatorname{Lie}(\mathrm{SO}(n))=\{X\in\mathfrak{gl}_n(\mathbb{R}),\ X+{}^{t}X=0\}$ 是反对称矩阵的集合.其维数为 $n(n-1)/2$.
\item $\mathfrak{so}_n(\mathbb{C})=\operatorname{Lie}(\mathrm{O}(n,\mathbb{C}))=\{X\in\mathfrak{gl}_n(\mathbb{C}),\ X+{}^{t}X=0\}$.其实维数为 $n(n-1)$.
\item $\mathfrak{u}_n=\operatorname{Lie}(\mathrm{U}(n))=\{X\in\mathfrak{gl}_n(\mathbb{C}),\ X+{}^{t}\overline{X}=0\}$.例如 $\mathfrak{u}_1=i\mathbb{R}\subset\mathfrak{gl}_1(\mathbb{C})=\mathbb{C}$.其维数为 $n^{2}$.
\item $\mathfrak{su}_n=\operatorname{Lie}(\mathrm{SU}(n))=\mathfrak{u}_n\cap\mathfrak{sl}_n(\mathbb{C})$.其维数为 $n^{2}-1$.
\item $\mathfrak{sp}_{2n}(K)=\operatorname{Lie}(\mathrm{Sp}_{2n}(K))=\{X\in\mathfrak{gl}_{2n}(K),\ {}^{t}XJ_{2n}+J_{2n}X=0\}$.为计算维数，注意 ${}^{t}XJ_{2n}+J_{2n}X=0\iff(J_{2n}X\ \text{对称})$.求得其实维数在 $K=\mathbb{R}$ 时为 $n(n+1)/2$，在 $K=\mathbb{C}$ 时为 $n(n+1)$.\footnote{译注：原文所列辛 Lie 代数的维数照录. 由其定义的矩阵方程分块计数，$K=\mathbb R$ 时实维应为 $n(2n+1)$，$K=\mathbb C$ 时实维应为 $2n(2n+1)$.}
\end{enumerate}

\begin{remark*}
定义 $\mathfrak{sl}_n$、$\mathfrak{so}_n$ 与 $\mathfrak{sp}_{2n}$ 的方程在\textbf{任意域} $K$ 上都有意义，从而给出一般域上 Lie 代数的例子.

此外，若给定 $K$ 的一个\newterm{ssl-quadratic-extension}{二次扩张}{quadratic extension} $K'$，还可以在 $\mathfrak{gl}_n(K')$ 中定义与 $\mathfrak{gl}_n(\mathbb{C})$ 中的 $\mathfrak{u}_n$ 类似的子 $K$-Lie 代数 $\mathfrak{u}_n(K'/K)$.
\end{remark*}
\end{numpara}

\begin{numpara}[Lie 群与其 Lie 代数之间的相互作用]\label{np:ssl-1-3-5}
下面是 Lie 群或 Lie 群同态的一些性质，它们可以从其 Lie 代数或诱导的 Lie 代数同态中读出.
\begin{itemize}
\item 连通 Lie 群 $G$ 的\textbf{连通} Lie 子群 $H$ 是\newterm{ssl-normal-subgroup}{正规子群}{normal subgroup}，当且仅当 $\operatorname{Lie}(H)$ 是 $\operatorname{Lie}(G)$ 的\newterm{ssl-ideal}{理想}{ideal}（即 $[\operatorname{Lie}(H),\operatorname{Lie}(G)]\subset\operatorname{Lie}(G)$\footnote{译注：原文此处包含于 $\operatorname{Lie}(G)$；按理想定义应为 $[\operatorname{Lie}(H),\operatorname{Lie}(G)]\subset\operatorname{Lie}(H)$ 之误.}）.
\item 连通 Lie 群 $G$ 的\textbf{连通} Lie 子群 $H$ 是\newterm{ssl-central-subgroup}{中心子群}{central subgroup}，当且仅当 $[\operatorname{Lie}(H),\operatorname{Lie}(G)]=0$.
\item 若 $G$ 连通，则伴随表示的核是 $G$ 的中心.
\end{itemize}
设 $\varphi\colon G\longrightarrow G'$ 是 Lie 群同态.
\begin{itemize}
\item $d\varphi$ 满射 $\iff\varphi(G)\supset G'^{0}$.
\item $\operatorname{Lie}(\operatorname{Ker}\varphi)=\operatorname{Ker}(d\varphi)$.特别地，$d\varphi$ 单射 $\iff\operatorname{Ker}(\varphi)$ 是离散的.
\item 若 $G'$ 连通，则 $d\varphi$ 双射 $\iff\varphi$ 是一个\newterm{ssl-covering}{覆盖}{covering}.
\end{itemize}
正如这些性质所示，当群连通时，从群过渡到其 Lie 代数更有意义（想一想：有限群的 Lie 代数为零……），而当群\newterm{ssl-simply-connected}{单连通}{simply connected}时意义还要更大.例如，[GALM2] 的定理 1.9.2 说：

\begin{theorem*}
设 $G$、$G'$ 是两个 Lie 群.映射
\[
\begin{array}{rcl}
\operatorname{Hom}_{Gp.Lie}(G,G') &\longrightarrow& \operatorname{Hom}_{\mathbb{R}-AL}(\operatorname{Lie}(G),\operatorname{Lie}(G'))\\
\varphi &\longmapsto& d\varphi
\end{array}
\]
是：
\begin{enumerate}[label=\roman*)]
\item 单射，若 $G$ 连通；
\item 双射，若 $G$ 单连通.
\end{enumerate}
\end{theorem*}

此外还可以证明：任何有限维 $\mathbb{R}$-Lie 代数都来自某个 Lie 群.于是有如下令人印象深刻的推论：函子 $G\mapsto\operatorname{Lie}(G)$ 是范畴等价
\[
\{\text{单连通 Lie 群}\}\ \xrightarrow{\ \sim\ }\ \{\text{有限维 }\mathbb{R}\text{-Lie 代数}\}.
\]
\end{numpara}

\begin{numpara}[复情形]\label{np:ssl-1-3-6}
若 $G$ 是复 Lie 群，则 $\operatorname{Lie}(G)$ 自然地是一个 $\mathbb{C}$-Lie 代数.此外，\newterm{ssl-holomorphic}{全纯}{holomorphic}同态 $\varphi\colon G\longrightarrow G'$ 的微分 $d\varphi$ 是 $\mathbb{C}$-线性的.
\end{numpara}
